\chapter{User Test Feedback Data}
\label{appendix:user_test_feedback}

This appendix summarizes the user-test feedback used in the evaluation of SafeVision. Instead of presenting each individual test form in full, the feedback was grouped into recurring themes. The tables show which features or interaction areas were commented on, the main feedback given by participants, and how the feedback influenced the system evaluation and design decisions. When it comes to Skan-Kontroll these user tests counts as multiple combined to one response.

\section{First User Testing Phase}
\label{app:user_test_phase_1}

The first user testing phase focused on the early version of the system. The
main purpose was to identify major usability issues related to workflow
understanding, navigation, terminology, and manual redaction. The feedback was
coded into the recurring themes below.

\subsection*{Workflow and navigation}

Several users were unsure what to do after uploading an image. In particular,
participants did not always understand that detection should be started before
redaction, and some users found it difficult to navigate from file upload to
detection and then to redaction.

\textbf{Participants:} ID001, ID002, ID003, ID004, ID006, ID009, Thomas.

\textbf{Design implication:} The workflow needed clearer step ordering, stronger
visual guidance after upload, and a button layout that better matched the
expected task progression.

\subsection*{Button naming and terminology}

The term ``AI mission'' was unclear to several users. Participants suggested
that the detection action should have a clearer and more intuitive name.

\textbf{Participants:} ID003, ID004, ID006, Thomas.

\textbf{Design implication:} The detection feature should use terminology that
more directly communicates its purpose, such as ``Identify'' or ``Detect''.

\subsection*{Manual redaction and editing tools}

Several users did not notice that manual redaction tools existed, or did not
fully understand that detected boxes could be edited. Some users also found it
unclear when manual tools were active.

\textbf{Participants:} ID001, ID002, ID005, ID007, ID008, Thomas.

\textbf{Design implication:} Manual editing tools needed clearer labels, more
visible interaction states, and better feedback when editing modes were active.

\subsection*{Settings and side panels}

Some users were unsure when and how detection or redaction settings should be
changed. The relationship between settings and the current workflow step was
not always clear.

\textbf{Participants:} ID001, ID003, ID005, ID009.

\textbf{Design implication:} Settings needed to be more visible and better
connected to the workflow step where they were relevant.

\subsection*{Progress and system feedback}

Users reported confusion around progress bars and system state, especially
when several files were involved or when processing took longer than expected.

\textbf{Participants:} ID007, ID008.

\textbf{Design implication:} Processing feedback should more clearly show what
the system is doing, which file is being processed, and whether the system is
still working.

\subsection*{Tooltips and onboarding}

Several users suggested tooltips, information boxes, or a short onboarding
sequence to explain buttons and guide first-time use.

\textbf{Participants:} ID002, ID003, ID006, ID009, Thomas.

\textbf{Design implication:} Lightweight guidance could reduce first-time
confusion without making the interface more complex.

\subsection*{Positive feedback on automation and simplicity}

Users generally reacted positively to automatic detection and redaction,
especially when faces were detected quickly. Several users also liked that the
system did not contain too many functions and became easier to use after the
first attempt.

\textbf{Participants:} ID001, ID002, ID003, ID004, ID006, ID007, ID008, ID009, Thomas.

\textbf{Design implication:} Automatic detection should remain central to the
system, while the interface should remain simple and focused.

\section{Second User Testing Phase}
\label{appendix:user_test_phase_2}

The second user testing phase focused on a more developed version of SafeVision.
At this stage, the system had undergone several improvements based on earlier
feedback, and users were generally able to complete the main workflow. The
feedback in this phase was therefore more focused on interaction details,
processing feedback, editing precision, file handling, and additional features
that could improve the practical use of the system.

\subsection*{Filter and settings visibility}

Users responded positively to being able to choose which filters should be
active during detection and redaction. However, some participants found the
filter and settings sections difficult to discover or did not immediately
understand when these options should be used.

\textbf{Participants:} ID010, ID011, ID012, ID014, Reell 001.

\textbf{Design implication:} Filter settings should be made more visible and
more clearly connected to the detection and redaction workflow.

\subsection*{Manual editing and tool state}

Several users commented on manual editing. Some users were unsure whether the
manual tool was active, especially when the active state was only indicated by
a highlighted button. Users also suggested more precise editing interactions,
such as deleting selected boxes with Backspace, right-click deletion, resizing
boxes, copying boxes, or drawing custom redaction areas.

\textbf{Participants:} ID008, ID009, ID010, ID011, ID012, ID014, ID015, ID016, Reell 001.

\textbf{Design implication:} Manual editing should provide clearer active-state
feedback and support more flexible interactions for editing, deleting, and
adjusting redaction regions.

\subsection*{Overlapping bounding boxes}

Some users found it difficult to delete or adjust a box when it overlapped with
another bounding box. In particular, it was unclear which box was selected or
which box would be affected by the next action.

\textbf{Participants:} ID014, Reell 001.

\textbf{Design implication:} The system should make selected regions clearer
and provide better support for editing overlapping detections.

\subsection*{Detection and redaction feedback}

Some participants reported that detection or redaction appeared to become stuck,
especially when progress stayed at 92\% for a longer period. This created
uncertainty about whether the system was still processing or had failed.

\textbf{Participants:} ID013, Reell 001.

\textbf{Design implication:} Progress indicators should provide more reliable
and informative feedback during longer processing tasks, especially when
processing multiple files or videos.

\subsection*{Detection quality and false positives}

Detection was generally perceived as useful, but some incorrect detections were
observed. For example, one test reported that shoelaces were detected as a
license plate. The domain stakeholder also emphasized that automatic detection
cannot be treated as fully reliable in all real-world cases.

\textbf{Participants:} Skan-Kontroll, Reell 001.

\textbf{Design implication:} Automated detection should remain combined with
manual review and correction, since users need to verify and adjust results
before final redaction.

\subsection*{Redaction modes and preview}

Users found the different redaction methods useful, especially the ability to
choose between different types of anonymization. Some users suggested that it
would be useful to preview the selected redaction mode before applying it, and
to adjust the strength or intensity of certain effects.

\textbf{Participants:} Skan-Kontroll, ID010, ID011, Reell 001.

\textbf{Design implication:} Redaction modes should provide clearer previews
and, where relevant, adjustable parameters such as intensity.

\subsection*{Metadata handling}

The domain stakeholder pointed out that metadata may also contain sensitive
information. It was therefore considered important that the system makes it
clear what happens to metadata during processing, and whether metadata is kept,
removed, or selectively handled.

\textbf{Participants:} Skan-Kontroll.

\textbf{Design implication:} Metadata handling should be made transparent and
treated as part of the privacy-preserving workflow.

\subsection*{Download behavior}

One participant suggested that downloaded files should automatically receive a
modified filename, for example by adding ``Anonymisert'' or ``Redacted'' to the
original filename. This would make it easier to distinguish the processed file
from the original.

\textbf{Participants:} Reell 001.

\textbf{Design implication:} Exported files should be clearly distinguishable
from the original input files.

\subsection*{Keyboard shortcuts}

A Mac user noted that shortcuts should work with the Command key rather than
requiring the user to manually change shortcuts from Control. This reflects
platform-specific expectations for keyboard interaction.

\textbf{Participants:} Reell 001.

\textbf{Design implication:} Keyboard shortcuts should follow platform
conventions where possible.

\subsection*{Video support and future workflow needs}

Video support was mentioned as important for realistic professional use cases.
The domain stakeholder and several users considered video redaction relevant,
especially in workflows involving larger amounts of visual material.

\textbf{Participants:} Skan-Kontroll, ID010, ID011.

\textbf{Design implication:} Video support is important for the practical value
of the system and should remain a central part of further development.

\subsection*{Overall impression}

Overall, users described the system as useful, professional, fast, and promising.
The final version was generally perceived positively, although users still
identified several smaller improvements related to feedback, editing precision,
and workflow clarity.

\textbf{Participants:} Skan-Kontroll, ID015, ID016, ID017, Reell 001.

\textbf{Design implication:} The system was perceived as practically useful, but
manual user control should remain central so that users can handle edge cases
and correct detection errors.

\subsection*{Final real user test with Skan-Kontroll}
\label{appsec:finale_real_test}

The final real user test with Skan-Kontroll provided feedback from a realistic professional use context. Overall, the feedback was highly positive. The system was described as impressive, intuitive, easy to navigate, and visually clean. The participants also stated that the solution could have been useful in their daily work if they had access to it after the project.

Several practical improvement areas were also identified. The participants noted that the patch notes were useful, but that their location or purpose in the interface was not fully clear. They also suggested that a simple FAQ or help page could support first-time users. Although the workflow was generally considered intuitive, some terminology was still seen as challenging. In particular, the word ``Redact'' was considered unfamiliar, and optional Norwegian language support was suggested as a possible improvement.

Processing feedback was another important issue. The participants observed that the system could take time when many filters were selected, and that the progress bar sometimes remained at 92\% for a long period. This made it appear as if the system had stopped, even when processing was still ongoing. They also noted that redaction did not provide the same level of progress feedback, and suggested that a clearer progress indicator or estimated completion time could reduce user uncertainty during longer tasks.

The participants also responded positively to the inclusion of video support. They considered video redaction to have strong commercial potential for future development, especially for organizations that handle larger amounts of visual material.

\textbf{Participants:} Skan-Kontroll.

\textbf{Design implication:} The feedback confirms that SafeVision has practical value in a professional setting, but also shows that future versions should improve user guidance, terminology, language support, and progress feedback. Video support should also remain an important direction for further development.